\documentclass[10pt,a4paper]{amsart}
\usepackage[margin=19mm]{geometry}
\usepackage{amsmath,amssymb,mathtools}
\usepackage[scr=rsfs,cal=euler]{mathalfa}
\usepackage{xfrac}
\usepackage[unicode,hidelinks,bookmarks=false]{hyperref}
\newcommand{\ie}{i.e.\ }
\newcommand{\cf}{cf.\ }
\newcommand{\resp}{respectively\ }
\newcommand{\Subsection}{\subsection}
\providecommand{\nobreakdash}{\nobreak\hbox{-}}
\setlength{\emergencystretch}{2em}
\multlinegap0pt
\usepackage{amssymb}
\usepackage[shortlabels]{enumitem}
\usepackage{xcolor}
\usepackage{tikz}
\usepackage[all]{xy}
\newtheorem{proposition}{Proposition}
\newcommand{\G}{\mathcal{G}}
\renewcommand{\P}{\mathcal{P}}
\newcommand{\init}{{\tt in}}
\newcommand{\qand}{\quad \mbox{and} \quad}
\newcommand{\qforall}{\quad \mbox{for all} \quad}
\newcommand{\qforeach}{\quad \mbox{for each} \quad}
\title{Polarization and Gorenstein Liaison}
\author{Sara Faridi, Patricia Klein, Jenna Rajchgot,, Alexandra Seceleanu}
\date{Source: arXiv:2406.19985v2 (CC BY 4.0)}
\begin{document}
\maketitle

\noindent\emph{Source and licence.} Polarization and Gorenstein Liaison. Version arXiv:2406.19985v2 (CC BY 4.0), distributed by arXiv under the Creative Commons Attribution 4.0 International licence, http://creativecommons.org/licenses/by/4.0/, as recorded in the arXiv licence metadata for this exact version and shown on the abstract page https://arxiv.org/abs/2406.19985v2. Full licence notice: \texttt{licenses/arxiv-2406.19985-CC-BY-4.0.txt}.\par\smallskip

\noindent\textit{Editorial scope.} Counted unit: the article's proposition prop:uniform-deg-in-y-polarizes (source0.tex lines 1060--1062). The article's complete original proof of it is delivered with it (source0.tex lines 1065--1120, one \texttt{proof} environment of 56 lines). No mathematics is written, repaired or re-derived: the statement and the proof are the article's own lines, in the article's own notation, and no retained line is substituted or re-flowed.  No further source-local context is needed: every cross-reference of every retained line resolves inside the two delivered spans.  Nothing was omitted from any retained span, so the package carries no omission record.  No retained line contains a citation, so the wrapper carries no bibliography and no bibliography entry.  No retained line contains a figure environment, an image-inclusion command or drawing code, so no image is delivered and the package declares no asset dependency.  Editorial formatting only, in addition: an AMS \texttt{amsart} preamble that restates the article's own declarations and macros used by the retained lines, namely the environments \texttt{proposition} on separate counters, together with the standard packages those lines need.  The retained environments print in this document as Proposition 1 in order of appearance, whereas the article prints the same items with the numbering of its own full document.  The complete native source of the article as distributed in the arXiv e-print for this exact version is bundled unchanged as \texttt{sources/161522/source0.tex}.
\begin{proposition}\label{prop:uniform-deg-in-y-polarizes}
Fix a homogeneous ideal $I$ of $R = \kappa[x_1, \ldots, x_n]$, a variable $y$ of $R$, a new variable $y'$, a $y$-compatible term order $<$, and a $<$-Gr\"obner basis $\G$ of $I$. Suppose there is some integer $t$ such that $\deg_y (g) \in \{0,t\}$ for each $g\in \G$. Then $\P_y (\G)$ is a Gr\"obner basis under the induced term order $\prec$ on $R' = R[y']$.
\end{proposition} 

\begin{proof}
Let $g,h\in \G$. Without loss of generality, assume that $g$ and $h$ are monic.  We will consider the $s$-polynomial $s(\P_y(g),\P_y(h))$ and show that its remainder upon division by $\P_y (\G)$ is $0$.  To do this, we will first fix a standard expression \[
s(g,h) = c_1g_1+\cdots+c_rg_r 
\] for $s(g,h)$ in terms of $\G$ and show that \[
\P_y(s(g, h)) = 
s(\P_y(g),\P_y(h)) 
\qand 
\P_y(c_i g_i) = 
\P_y(c_i)\P_y(g_i) \qforall i \in [r],
\] from which it follows that \[ s(\P_y(g),\P_y(h))=\P_y(c_1)\P_y(g_1)+\cdots+\P_y(c_r)\P_y(g_r).
\]

In order to establish this claim, we will frequently make use of the observation that, if $u, v \in R$ and $\deg_y(u)=0$, then $\P_y(uv) = u \P_y(v) = \P_y(u)\P_y(v)$.  We consider three cases.

\noindent \textbf{Case 1:} Suppose first that $\deg_y(g) = \deg_y(h) = 0$. Then $\deg_y(c_ig_i) = 0$ for all $i \in [r]$, and so the claim follows from the above observation.



\noindent \textbf{Case 2:} Next suppose $\deg_y(g) = \deg_y(h) = t$. 
%Then we may write $h = y'd_1+r_1$
Then there are some monomials $m$ and $n$ such that 
$s(\P_y(g),\P_y(h)) = m\P_y(h) - n\P_y(g)$. 
By the assumption that $\deg_y(g) = \deg_y(h) = t$, we observe that neither $m$ nor $n$ is divisible by $y$ or $y'$. 
Consequently, 
\[
    s(\P_y(g),\P_y(h)) = m\P_y(h) - n\P_y(g)
     = \P_y(mh) - \P_y(ng)   = \P_y(s(g,h)).
\]
Now recall the standard expression
$s(g,h) = c_1g_1+\cdots + c_rg_r$ for $s(g,h)$ in terms of $\G$.
For each $i \in [r]$, if $g_i$ has a term divisible by $y$, then, by assumption, $\deg_y(g_i) = t$. Hence, because $\deg_y(c_ig_i) \leq \deg_y (s(g,h))\leq t$, we must have that $\deg_y(c_i) = 0$. 
Hence, for all $i \in [r]$, $\P_y(c_ig_i) = \P_y(c_i)\P_y(g_i)$ by our earlier observation. 

\noindent \textbf{Case 3:} Without loss of generality, assume $\deg_y(g) = t$ and $\deg_y(h) = 0$.
    Then there are monomials $m$ and $n$ such that 
$s(\P_y(g),\P_y(h)) = my(y')^{t-1}\P_y(h) - n\P_y(g)$,
where neither $m$ nor $n$ is divisible by $y$ or $y'$. Consequently, 
\[
    s(\P_y(g),\P_y(h)) = s(\P_y(g),h)= my'y^{t-1}h - n\P_y(g)
     = \P_y(my^t h) - \P_y(ng)   = \P_y(s(g,h)).
\]
The remainder of the claim in this third and final case follows via the argument from Case~2.

Having established that, for any $g, h \in \G$ and any standard expression \[
s(g,h) = c_1g_1+\cdots+c_rg_r,
\] we have \[
 s(\P_y(g),\P_y(h)) = \P_y(c_1)\P_y(g_1)+\cdots+\P_y(c_r)\P_y(g_r),
\] it remains only to show that \[
\init_\prec(s(\P_y(g),\P_y(h))) \geq \init_\prec(\P_y(c_i)\P_y(g_i)) \qforeach  i \in [r].
\] From the equality $s(\P_y(g),\P_y(h)) = \P_y(s(g, h))$, it follows that \[
\init_\prec(s(\P_y(g),\P_y(h))) = \init_\prec(\P_y(s(g, h))) = \P_y(\init_<(s(g,h))).
\] Similarly, from the equalities $\P_y(c_i)\P_y(g_i) = \P_y(c_ig_i)$ for each $i \in [r]$, it follows that \[
\init_\prec(\P_y(c_i)\P_y(g_i)) = \init_\prec(\P_y(c_ig_i)) = \P_y(\init_<(c_ig_i)).
\]
The assumption $\init_<(s(g,h)) \geq \init_<(c_ig_i)$ implies $\P_y(\init_<(s(g,h))) \geq \P_y(\init_<(c_ig_i))$.  Combining this information, we see $\init_\prec(s(\P_y(g),\P_y(h))) \geq \init_\prec(\P_y(c_i)\P_y(g_i))$ for all $i \in [r]$, which shows that $s(\P_y(g),\P_y(h))$ has remainder $0$ upon division by $\P_y(G)$.  Hence $\P_y(\G)$ forms a Gr\"obner basis.
\end{proof}

\end{document}
